\section{Implementation}
\label{sec:implementation-ra}

This checked-in artifact does not ship an endpoint enforcement engine.
It specifies the implementation boundary an endpoint engine must satisfy:
twelve action variants, four reversible variants with inverse executors,
content-hash-gated rollback constructors, and a four-receipt grammar. The
substrate inherits the canonical-byte signing path, ledger discipline,
and verifier-owned trust store conventions from the parent paper. Claims
about response executors and proof packages below are implementation
obligations until their code lands in this repository.

\paragraph{Action plan and execution.}
The response plan is bound at request time:
\codepath{EndpointResponsePlan} carries action identifier, action
class, dry-run flag, root node identifier, graph slice identifier,
TTL in seconds, rollback reference handle, reason, issued-at and
expires-at timestamps, and node and edge counts. The expires-at
timestamp is derived structurally from issued-at and TTL. The TTL
validator rejects zero and out-of-range durations at admission; the
upper bound is the deployment-instance constant
\codepath{EDR_MAX_RESPONSE_TTL_SECONDS}.

The forward executor dispatches on the action class. File quarantine
is a hashed read followed by \codepath{fs::rename} into the quarantine
root. Persistence disable is the same shape with a different
destination subtree under the same quarantine root. Process-tree
suspend issues \codepath{libc::SIGSTOP} to each pid in the target
list in reverse-order traversal of the process tree. Egress
restriction writes a policy snapshot file at the network-extension
policy path and invokes the network-extension reload entry point. Each
executor returns a signed execution receipt that records the forward
effect set; an effect carries effect identifier, effect type, target,
artifact path, content hash, and byte count.

\paragraph{Rollback executor.}
The rollback executor's inverse pair runs on the four reversible
variants. File quarantine's inverse reads the artifact, hashes it,
verifies content-hash parity against the recorded forward hash,
refuses the operation if the original target path now exists, and
issues the inverse \codepath{fs::rename}. Persistence disable's
inverse has the same shape. Process-tree suspend's inverse issues
\codepath{libc::SIGCONT} to each pid; the substrate records which pids
in the original set have since exited and reports the partial-inverse
status. Egress restriction's inverse deactivates the egress ledger
entry and re-syncs the network-extension policy. Each inverse
executor emits a rollback receipt
(\codepath{EndpointResponseRollbackReport}) that references the
execution identifier and the action identifier, and the receipt is
appended to the rollback ledger before the next dispatch.

\paragraph{Acknowledgement loop.}
The acknowledgement receipt
(\codepath{EndpointResponseAcknowledgementReport}) is emitted by an
endpoint-initiated postback to the control plane. The acknowledgement
identifier, execution identifier, action identifier, acknowledged-by
key, and optional control-correlation token are bound to the receipt.
The substrate does not mutate execution-ledger state on
acknowledgement; the acknowledgement is an additive receipt event the
verifier-owned correlator ingests.

\paragraph{Receipt signing and verification.}
Each response receipt is signed by the endpoint's Ed25519 keypair over
the canonical-JSON body. The signing path is shared across the four
receipt families, which is the substrate's chosen trust granularity
for the response engine. Bilateral cosignature on the destructive
subclass is a substrate-level obligation discussed in
Section~\ref{sec:limitations-ra}: the deployment instance carries the
single-signer path; the bilateral path is a substrate primitive
inherited from the parent paper's
\codepath{crates/trust/chio-federation/src/bilateral_dsse.rs} that the
response engine is not wired against.

\paragraph{Proof bundle.}
No checked-in proof endpoint backs this paper yet. The required proof
bundle shape is: execution record, graph slice, provider and sensor
state, and six receipt families (request, execution,
evidence-bundle-manifest, lifecycle transitions, rollback,
acknowledgement). A future endpoint
\codepath{GET /api/v1/agent/edr/response-executions/\{id\}/proof}
must return that bundle and let the verifier replay canonical bytes
against the polity's pinned public keys while asserting actor continuity
across the chain. Until then, the proof bundle is a design contract, not
a deployed implementation claim.
